% Lösungen Einheit 1 von 4 – Normalparabel und Streckfaktor (zusammenbau v0.1)
\einheitenkopf{Einheit 1 von 4 · Normalparabel und Streckfaktor}

%% TODO Lösung der Erklärzeile 11g) fehlt in der Bank
\erg{11}{a) Gerade \quad b) $4$; $1$; $0$; $1$; $4$ \quad c) $9$; $4$; $1$; $0$; $1$; $4$; $9$ \quad d) $16$; $4$; $0$; $4$; $16$ \quad e) $25$; $9$; $1$; $1$; $9$ \quad f) $36$; $16$; $1$; $9$; $25$ \quad g) \ldots \quad h) Bogen ohne Lineal durch $(-3|9)$, $(-2|4)$, $(-1|1)$, $(0|0)$, $(1|1)$, $(2|4)$, $(3|9)$; tiefster Punkt $(0|0)$ \quad i) $S(0|0)$; Symmetrieachse ist die $y$-Achse \quad j) $f(-6) = (-6)^2 = 36$ \quad k) $16$; $4$; $0$; $4$; $16$ \quad l) nach unten, schmaler als die Normalparabel \quad m) Tabelle A: $0{,}5 \cdot 2^2 = 2$ \quad n) Bogen durch $(-4|8)$, $(-2|2)$, $(0|0)$, $(2|2)$, $(4|8)$ \quad o) $f(x) = -x^2 + 3$: nach unten geöffnet, $f(0) = 3$ \quad p) Tabelle B: $4 \cdot 2^2 = 16$ – erst quadrieren, dann malnehmen \quad q) $h(t) = -5t^2 + 45$: Startwert $h(0) = 45$, der Bogen fällt, also Minus vor $t^2$; Kontrolle $h(2) = 25$ passt zum Graphen.}
\erg{12}{a) Negative $x$ negativ quadriert; richtig: $(-4)^2 = 16$, $(-2)^2 = 4$ – Tabelle $16$; $4$; $0$; $4$; $16$ \quad b) Beim Quadrieren fällt das Vorzeichen weg: $(-8) \cdot (-8) = 8 \cdot 8 = 64$ – die Parabel ist symmetrisch zur $y$-Achse. \quad c) $a = 7$, denn $f(1) = a \cdot 1^2 = 7$; $f(x) = 7x^2$ \quad d) $h(20) = -0{,}02 \cdot 400 = -8$; der Bogen liegt dort $8\,$m tiefer.}
